\section{Evaluation}
\label{sec:evaluation}

This section compares the four classical segmentation methods on the
parametric staircase point cloud. The setup and metrics come first (A),
followed by the quantitative results and their interpretation (B). We then
examine robustness as height noise grows (C) and close with the safety
implications for stair-following robots (D).

\subsection{Experimental Setup and Metrics}

All experiments ran on the synthetic staircase from
Section~\ref{sec:method}, with $N=4$ steps, width $w=2.5$\,m, tread depth
$d=0.30$\,m, and step height $h=0.18$\,m. The generated cloud holds
$18{,}333$ points: $12{,}800$ tread points (label~0), $4{,}800$ riser
points (label~1), and $733$ ``other'' points (label~2) that come from
edges, LiDAR jitter, and the injected outliers.
Figure~\ref{fig:profile} shows the $y$--$z$ side profile of the
ground-truth cloud. The alternating horizontal treads and vertical risers
stand out clearly, along with the sparse off-surface clutter.

\begin{figure}[t]\centering
\includegraphics[width=0.95\columnwidth]{profile_ground_truth.png}
\caption{Side profile ($y$--$z$) of the ground-truth cloud: green = tread, red = riser, grey = other.}
\label{fig:profile}
\end{figure}

The steppable surface is the safety-critical class, so we score every
method on the tread label in a one-vs-rest setting: tread points are
positives, and all riser and ``other'' points are negatives.
Let $TP$, $FP$, $FN$, and $TN$ denote the per-point true positives, false
positives, false negatives, and true negatives. We report five standard
quantities, each computed for the tread class under this one-vs-rest
scheme. The tabulated accuracy $\mathrm{Acc}=(TP{+}TN)/(TP{+}TN{+}FP{+}FN)$
is the fraction of points placed correctly inside or outside the tread
class. In the text we also quote the overall \emph{three-class} accuracy,
the fraction of all points whose full tread/riser/other label is correct;
this differs from the one-vs-rest $\mathrm{Acc}$ of
Table~\ref{tab:metrics}. Precision
$\mathrm{P}=TP/(TP{+}FP)$ is the fraction of predicted tread points that
are genuinely treads, while recall $\mathrm{R}=TP/(TP{+}FN)$ is the
fraction of true tread points recovered. The F1 score is their harmonic
mean, $\mathrm{F1}=2\mathrm{P}\mathrm{R}/(\mathrm{P}{+}\mathrm{R})$, and the
intersection-over-union, or Jaccard index,
$\mathrm{IoU}=TP/(TP{+}FP{+}FN)$, expresses the spatial overlap between the
predicted and ground-truth tread regions. Precision and IoU matter most
for robotic foothold reasoning, because they directly measure how often a
non-steppable surface is taken for a steppable one.

Reproducibility was a design goal. The data generator and every stochastic
component of the pipeline (RANSAC sampling, neighbour selection, and the
noise draws) ran from a fixed pseudo-random seed, so the cloud and all
reported numbers regenerate deterministically. Method hyper-parameters
stayed constant across runs: RANSAC used an inlier threshold
$\tau=0.035$\,m, PCA normals were estimated from $k=50$ neighbours, and
DBSCAN used $\varepsilon=0.08$\,m within each orientation group.

\subsection{Results}

Table~\ref{tab:metrics} summarises the tread-class performance of the four
methods, and Fig.~\ref{fig:bars} places their F1 and IoU scores side by
side.

\begin{table}[t]
\caption{Tread-class segmentation metrics (one-vs-rest). Acc is one-vs-rest
tread accuracy, distinct from the three-class accuracy quoted in the text.
Best per column in bold.}
\label{tab:metrics}
\centering
\begin{tabular}{lccccc}
\toprule
Method & Acc & Prec & Rec & F1 & IoU \\
\midrule
RANSAC          & 0.681 & 0.714 & 0.904 & 0.798 & 0.664 \\
PCA-Normals     & 0.901 & \textbf{0.954} & 0.901 & 0.927 & 0.864 \\
DBSCAN-Slope    & \textbf{0.920} & 0.936 & 0.949 & \textbf{0.943} & \textbf{0.892} \\
Height-Histogram& 0.827 & 0.809 & \textbf{0.984} & 0.888 & 0.799 \\
\bottomrule
\end{tabular}
\end{table}

\begin{figure}[t]\centering
\begin{tikzpicture}
\begin{axis}[width=\columnwidth,height=4.6cm,ybar,bar width=6pt,ymin=0,ymax=1,
  symbolic x coords={RANSAC,PCANormals,DBSCANSlope,HeightHistogram},
  xtick=data,xticklabel style={rotate=20,anchor=east,font=\footnotesize},
  ylabel={score},legend pos=south east,enlarge x limits=0.15]
\addplot table[x=method,y=f1]{data/metrics.dat}; \addlegendentry{F1}
\addplot table[x=method,y=iou]{data/metrics.dat}; \addlegendentry{IoU}
\end{axis}
\end{tikzpicture}
\caption{Per-method tread-class F1 and IoU.}
\label{fig:bars}
\end{figure}

The orientation-aware \textbf{DBSCAN-Slope} method leads overall. It
attains the best F1 ($0.943$) and IoU ($0.892$) and the highest three-class
accuracy ($0.909$). The reason is its order of operations: by first
splitting the cloud into horizontal and vertical orientation groups and
clustering spatially only afterwards, it recovers each physical surface as
a single compact cluster and absorbs outliers into the noise label. The
result is a balanced precision ($0.936$) and recall ($0.949$).
Figure~\ref{fig:seg3d} sets the DBSCAN-Slope segmentation against the
ground truth; the individual treads separate cleanly, and riser
bleed-through is minimal.

\begin{figure}[t]\centering
\includegraphics[width=\columnwidth]{seg_dbscan_slope.png}
\caption{3D segmentation by DBSCAN with slope analysis (right) vs.\ ground truth (left).}
\label{fig:seg3d}
\end{figure}

The \textbf{PCA-Normals} method follows close behind, with F1 $0.927$ and
IoU $0.864$. It classifies points directly from the local surface
orientation $|\nvec^\top\zhat|$, and this earns it the highest precision of
any method ($0.954$): a point is labelled tread only when its estimated
normal is nearly vertical, so very few risers slip through. Recall
($0.901$) sits marginally below DBSCAN's, because normals estimated near
tread--riser edges grow unreliable and a fraction of genuine tread points
fall outside the horizontal gate. The narrow gap between PCA-Normals and
DBSCAN-Slope confirms that local geometric orientation is an effective,
robust cue for separating the two principal surface types.

The \textbf{Height-Histogram} method acts as a high-recall,
lower-precision detector. It reaches the best recall of any method
($0.984$) but a markedly lower precision ($0.809$), giving F1 $0.888$ and
IoU $0.799$. Figure~\ref{fig:hist} shows the $z$-histogram, whose four
pronounced peaks correspond to the detected tread levels at
$z\approx\{0.00,0.18,0.36,0.54\}$\,m, close to the nominal step heights.
The method captures almost every tread point because treads concentrate
sharply at these levels. The trouble is that the horizontal bands it
assigns around each peak also enclose riser points whose vertical extent
happens to pass through a tread height, and that drags precision down.

\begin{figure}[t]\centering
\begin{tikzpicture}
\begin{axis}[width=\columnwidth,height=4.2cm,xlabel={$z$ (m)},ylabel={count},
  ymin=0,ymax=760,xmin=-0.1,xmax=0.65,enlarge x limits=false,tick label style={font=\footnotesize}]
\addplot[draw=blue!60!black,fill=blue!25] table[x=z,y=count]{data/hist.dat} \closedcycle;
\draw[green!55!black,dashed,thick] (axis cs:0.004,0) -- (axis cs:0.004,760);
\draw[green!55!black,dashed,thick] (axis cs:0.183,0) -- (axis cs:0.183,760);
\draw[green!55!black,dashed,thick] (axis cs:0.362,0) -- (axis cs:0.362,760);
\draw[green!55!black,dashed,thick] (axis cs:0.541,0) -- (axis cs:0.541,760);
\end{axis}
\end{tikzpicture}
\caption{Height histogram of $z$; dashed lines mark the four detected tread levels.}
\label{fig:hist}
\end{figure}

Plain \textbf{RANSAC} plane fitting is the weakest performer. It posts the
lowest precision ($0.714$), F1 ($0.798$), IoU ($0.664$), and three-class
accuracy ($0.667$), even with a high recall ($0.904$). One failure mode
dominates: the large horizontal planes RANSAC fits at each tread height
also swallow riser points that lie within the inlier band $\tau$ at those
heights. These mislabelled risers inflate the false-positive count and
collapse precision. The contrast with PCA-Normals is telling. A global
planar fit cannot tell a riser point that merely sits near a tread plane
from a true tread point, whereas a local normal can. This is why the
per-point orientation cues in PCA-Normals and DBSCAN-Slope beat the purely
planar formulation.

\subsection{Noise Sensitivity}

To probe robustness, we swept the height-noise standard deviation
$\sigma_z$ from $0.01$\,m to $0.07$\,m and re-evaluated PCA-Normals, whose
accuracy hinges most directly on the fidelity of local geometry. We vary
$\sigma_z$ rather than the riser depth jitter $\sigma_y$ because vertical
jitter is the axis that erodes the tread/riser orientation contrast,
whereas $\sigma_y$ mostly broadens riser planes without moving tread points
across step levels. Figure~\ref{fig:noise} plots the resulting tread-class
F1 and IoU.

\begin{figure}[t]\centering
\begin{tikzpicture}
\begin{axis}[width=\columnwidth,height=4.4cm,xlabel={$\sigma_z$ (m)},
  ylabel={tread-class score},ymin=0,ymax=1,legend pos=south west,
  grid=major,tick label style={font=\footnotesize}]
\addplot[blue,mark=*] table[x=sigma,y=f1]{data/noise.dat}; \addlegendentry{F1}
\addplot[red,mark=square*] table[x=sigma,y=iou]{data/noise.dat}; \addlegendentry{IoU}
\end{axis}
\end{tikzpicture}
\caption{Tread-class F1 and IoU versus height noise $\sigma_z$ (PCA-Normals).}
\label{fig:noise}
\end{figure}

Performance degrades monotonically, and the slope steepens as the noise
grows. At $\sigma_z=0.01$\,m the method is near-ideal (F1 $0.956$, IoU
$0.915$), and it holds up well through $\sigma_z=0.02$--$0.03$\,m. Past that
point the decline accelerates, reaching F1 $0.543$ and IoU $0.373$ at
$\sigma_z=0.07$\,m. The mechanism is geometric. The discriminative cue is
the near-vertical tread normal, but once $\sigma_z$ approaches half the step
height ($h/2=0.09$\,m), vertical jitter on tread points becomes comparable
to the vertical gap between adjacent steps. The local neighbourhoods used
for normal estimation then mix points from neighbouring treads and the
intervening riser, the estimated normals tilt away from vertical, the
horizontal orientation gate misfires, and tread and riser turn
geometrically indistinguishable. This marks a practical operating limit:
reliable tread segmentation needs the sensor's per-point height noise to
stay well below the step height it must resolve.

\subsection{Limitations and Threats to Validity}

Several limitations bound these conclusions. The benchmark is synthetic, a
single straight-flight, rectangular staircase with one fixed geometry
$(w,d,h,N)$ and an idealised, axis- and gravity-aligned layout. The
absolute scores should therefore be read as \emph{relative} method
rankings, not field-deployable accuracies. The sensor model is also
simplified: nominal Gaussian jitter plus uniform outliers, with no
occlusion, missing returns, incidence-angle effects, multi-path, or the
structured scan patterns of a real LiDAR. The height-histogram method, for
its part, assumes a regular, gravity-aligned staircase and would not
transfer to spiral, open-riser, or unaligned scans without modification.
Finally, every hyper-parameter ($\tau$, $k$, $\varepsilon$, the orientation
gates, and the histogram window) was tuned once on this cloud, and the
sensitivity study varies only $\sigma_z$ for a single method, so we make no
claim of cross-configuration parameter robustness. These constraints
motivate the real-sensor validation identified in
Section~\ref{sec:conclusion}.

\subsection{Safety Discussion for Stair-Following Robots}

For a stair-climbing or stair-following robot, the worth of these results
lies less in the average score than in the failure modes they expose. The
best method reaches IoU $\approx 0.89$, enough to support \emph{coarse}
foothold planning: the extracted tread regions overlap the true steppable
surfaces strongly enough for an elevation-map-based planner to place
candidate footholds on the correct steps, in line with practice in legged
and wheeled terrain perception
\cite{chavezgarcia2017learning,gomes2023survey}. Yet no method reaches unit
precision, and the dominant errors are exactly the dangerous ones: riser
points occasionally read as steppable. Inside a closed control loop such
false positives can lead a controller to commit a foothold to a vertical
face, the very failure mode that motivates dedicated stair perception
\cite{matsumura2022deep}.

We therefore recommend that a deployed system treat the raw per-point
labels as a hypothesis rather than a command. Three mitigations follow
directly from our findings. The first is \emph{ensemble voting} across
methods. Because the methods fail in complementary ways (RANSAC and the
histogram over-include risers, while PCA-Normals stays conservative),
voting can suppress method-specific false positives and combine the
histogram's high recall with the precision of the normal-based gate.
The second is a \emph{tread-edge erosion margin}. Mislabelling concentrates
at tread--riser boundaries, so shrinking each accepted tread region inward
before foothold selection trades a small amount of usable surface for a
sharp drop in the chance of stepping onto an edge or riser.
The third is \emph{temporal verification}. Given the strong noise
sensitivity established above, accumulating evidence across consecutive
scans before committing a foothold guards against transient mislabels
caused by momentary high-noise frames. Together these measures convert a
segmentation that is accurate on average into one whose residual errors
lean toward the safe side, which is the right design objective for robots
that physically traverse stairs.
